\documentclass[11pt]{article}
\usepackage[margin=1in]{geometry}
\usepackage{amsmath,amssymb,amsthm,mathtools}
\usepackage[T1]{fontenc}
\usepackage[expansion=false]{microtype}
\usepackage[hidelinks]{hyperref}
\newtheorem{theorem}{Theorem}[section]
\newtheorem{lemma}[theorem]{Lemma}
\newtheorem{proposition}[theorem]{Proposition}
\newcommand{\ddc}{\mathrm d\mathrm d^c}
\newcommand{\PSH}{\operatorname{PSH}}
\newcommand{\Vol}{\operatorname{Vol}}
\newcommand{\MA}{\operatorname{MA}}
\newcommand{\CP}{\mathbb P}
\newcommand{\norm}[1]{\left\|#1\right\|}
\title{Hausdorff stability of relative normalization functionals}
\author{}
\date{}
\begin{document}
\maketitle
\begin{abstract}
A common translated simplex gives an explicit geometric criterion for
stability of ordinary relative Monge--Amp\`ere normalization functionals
on projective space and a fixed point blowup. Hausdorff convergence of
the defining convex bodies gives a quantitative uniform bound on the
functionals over all continuous tests. The limiting reference is an
explicit admissible ordinary model, and its bounded-type solvability is
verified before identifying the limiting functional. Every supporting
probability is the same normalized exponential density. One concrete
family has isolated singularities that are not bounded-equivalent to
analytic singularities, including a positive-mass endpoint. A second
family has smooth residual forms of rank one on an exceptional divisor
and positive top mass. The result concerns these closed projective
backgrounds; unrestricted Hermitian degeneration remains open.
\end{abstract}

\section{The geometric criterion}

Use $\ddc=(i/(2\pi))\partial\bar\partial$, with exterior powers without
divided factorials. All potentials have their canonical upper
representatives. Write $v\preceq u$ if $v\le u+C$ and $v\sim u$ for
mutual bounded difference, including common poles. The ordinary envelope is
\[
 P_\theta[u]=\left(\lim_{C\to\infty}
 P_\theta(\min\{u+C,0\})\right)^*.
\]
An ordinary model reference satisfies $\sup_X\phi=0$ and
$P_\theta[\phi]=\phi$. We also verify a genuine
$P_{\varepsilon\theta}(\phi)$ for some $0<\varepsilon<1$.

Let $n\ge2$ and put
\[
 \Sigma=\{p\in\mathbb R^n:p_i\ge0,\ \textstyle\sum_i p_i\le1\}.
\]
On $T=(\mathbb C^*)^n\subset\CP^n$ use
\[
 x_i=\log|z_i|^2,\qquad L(x)=\log(1+\textstyle\sum_i e^{x_i}),
 \qquad\omega=\ddc\log(1+\textstyle\sum_i|z_i|^2).
\]
Thus $\int_{\CP^n}\omega^n=1$. Set
\[
 \ell(p)=\sum_i p_i\log p_i+
       (1-\textstyle\sum_i p_i)\log(1-\textstyle\sum_i p_i),
 \qquad 0\log0=0.
\]
For a convex body $K\subset\Sigma$ define
\[
 h_K(x)=\sup_{p\in K}p\cdot x,\qquad
 g_K(x)=\sup_{p\in K}\{p\cdot x-\ell(p)\}.
\]
Let $\psi_K$ and $\Phi_K$ denote the canonical extensions of
$h_K-L$ and $g_K-L$; their existence will be proved below.

Take $X=\CP^n$ with $\pi=\mathrm{id}$, or
$X=\operatorname{Bl}_q\CP^n$ for a fixed point $q$, with blowdown $\pi$.
Put $\Theta=\pi^*\omega$ and fix any smooth positive probability
volume $\mathrm dV$ on $X$. For an admissible ordinary reference write
$B_{\theta,\phi}(F)=\log c_{\theta,\phi}(F)$, where
\[
 \MA_\theta^{\mathrm{NP}}(u)=c_{\theta,\phi}(F)e^F\mathrm dV,
 \qquad u\sim\phi.
\]
Scalar uniqueness is part of the specified admissible framework.
Its supporting set $A_{\theta,\phi}(F)$ consists of probabilities $\nu$ with
\[
 B_{\theta,\phi}(G)\le B_{\theta,\phi}(F)
                   -\int_X(G-F)\,\mathrm d\nu
 \quad(G\in C(X,\mathbb R)).
\]

\begin{theorem}\label{thm:criterion}
Suppose convex bodies $P_j,P\subset\Sigma$ contain the same translated
simplex $p_*+\delta\Sigma$, where $0<\delta<1$, and
$h_j=d_H(P_j,P)\to0$ in Euclidean Hausdorff distance. Let
$r_j>0$ tend to $r>0$. Define
\[
 \theta_j=r_j\Theta,\quad \theta=r\Theta,\qquad
 \phi_j=r_j\pi^*\Phi_{P_j},\quad \phi=r\pi^*\Phi_P.
\]
The backgrounds are smooth closed semipositive and big, and all references,
including the limit, are ordinary models with genuine smaller-background
envelopes for $\varepsilon=1-\delta$. Their reference masses are
\[
 m_j=r_j^n n!\Vol(P_j),\qquad m=r^n n!\Vol(P)>0.
\]
Every bounded real Borel $F$ admits an ordinary-type solution at every
datum, including the limit, and with $Z(F)=\int_Xe^F\mathrm dV$ one has
\[
 c_j(F)=\frac{m_j}{Z(F)},\quad B_j(F)=\log m_j-\log Z(F),\qquad
 c(F)=\frac m{Z(F)},\quad B(F)=\log m-\log Z(F).
\]
Set $s=\delta/(n+\sqrt n)$. Then
\begin{align}
 \sup_{G\in C(X,\mathbb R)}|B_j(G)-B(G)|
 &=\left|\log\frac{r_j^n\Vol(P_j)}{r^n\Vol(P)}\right|\notag\\
 &\le n|\log(r_j/r)|+n\log(1+h_j/s).\label{eq:functional-bound}
\end{align}
Moreover $n!\Vol(P_j),n!\Vol(P)\ge\delta^n$,
$\theta_j\to\theta$ smoothly, and $\phi_j\to\phi$ in
$L^1(X,\mathrm dV)$ and locally uniformly on $\pi^{-1}(T)$.
For every continuous real $F$,
\[
 A_j(F)=A(F)=\{\nu_F\},\qquad
 \nu_F=\frac{e^F\mathrm dV}{Z(F)}.
\]
If $F_j\to F$ uniformly and $\nu_j\in A_j(F_j)$, then
\[
 \int_X\left|\frac{\mathrm d\nu_j}{\mathrm dV}
             -\frac{\mathrm d\nu_F}{\mathrm dV}\right|\mathrm dV
 \le e^{2\norm{F_j-F}_\infty}-1.
\]
Thus all support clusters are identified, with convergence in total
variation for uniformly varying tests.
\end{theorem}

The condition is geometric and supplies the limiting solvable class.
It does not assume functional convergence. The existence input is the
prescribed-singularity theorem of Darvas--Di Nezza--Lu
\cite[Theorem A(i); Theorems 4.3 and 4.7]{DDL}.
The convex and pluripotential dictionary is classical
\cite[Sections 2.1--2.3]{BB}\cite[Section 7]{DDL}; the explicit envelope, constants, common
admissibility witness and concrete families are established here.

\section{Canonical models and smaller-background witnesses}

Convex conjugation gives $L^*=\ell$ on $\Sigma$ and $+\infty$ elsewhere.
In the interior, $\nabla L(x)=p$ is equivalent to
$e^{x_i}=p_i/(1-\sum p_i)$; boundary values follow by limits. Since
$-\log(n+1)\le\ell\le0$, one obtains
\begin{equation}\label{eq:bounded-type}
 h_K\le g_K\le h_K+\log(n+1),\qquad g_K\le L,
 \qquad h_K\le h_\Sigma\le L.
\end{equation}
Convex functions of logarithmic moduli are psh on $T$: the assertion
follows from the Hessian for smooth convex functions and from convolution
for arbitrary finite convex functions. After adding a smooth local
projective potential, the upper bound by zero permits canonical psh
extension across the coordinate boundary
\cite[I.5.23--5.24]{Demailly}. Hence $\psi_K,\Phi_K$ are genuine
$\omega$-psh potentials and
$0\le\Phi_K-\psi_K\le\log(n+1)$ off their common poles.
Choose $p\in\operatorname{int}K\subset\operatorname{int}\Sigma$ and
$x$ with $\nabla L(x)=p$. Fenchel equality gives $g_K(x)=L(x)$, so
$\sup\Phi_K=0$.

\begin{lemma}\label{lem:model}
One has $P_\omega[\psi_K]=\Phi_K=P_\omega[\Phi_K]$.
\end{lemma}
\begin{proof}
For $C\ge\log(n+1)$ let
$v_C=P_\omega(\min\{\psi_K+C,0\})$. The reference $\Phi_K$ itself is
a competitor by \eqref{eq:bounded-type}. Canonical envelope regularization
\cite[I.5.7]{Demailly} gives $\Phi_K\le v_C\le0$ and
$v_C\le\psi_K+C$. Rotations preserve the background and obstacle,
so the envelope is invariant. On $T$ its local potential $f=v_C+L$ is
a finite convex function, with $f\le L$ and $f\le h_K+C$.
Separation shows $\operatorname{dom}f^*\subset K$; the first bound gives
$f^*\ge\ell$. Therefore
\[
 f(x)=\sup_{p\in\operatorname{dom}f^*}\{p\cdot x-f^*(p)\}
       \le g_K(x).
\]
Thus $v_C\le\Phi_K$ on $T$, and canonically at the boundary.
The opposite inequality was supplied by the actual competitor.
The finite envelopes therefore equal $\Phi_K$ once $C\ge\log(n+1)$.
Bounded-equivalent obstacles have interlacing finite-$C$ families;
hence their limiting envelopes agree. This proves both identities.
\end{proof}

The inclusion $p_*+\delta\Sigma\subset K\subset\Sigma$ implies
$p_{*,i}\ge0$ and $b=\sum p_{*,i}\le1-\delta=\varepsilon$.
On $T$ set
\[
 \chi=p_*\cdot x-\varepsilon L.
\]
Its canonical projective extension satisfies the exact current identity
\[
 \varepsilon\omega+\ddc\chi
 =\sum_i p_{*,i}[Z_i=0]+(\varepsilon-b)[Z_0=0]\ge0.
\]
Since $h_K\ge p_*\cdot x+\delta h_\Sigma$ and
$0\le L-h_\Sigma\le\log(n+1)$,
\[
 \Phi_K\ge\chi-\delta\log(n+1).
\]
This supplies an actual genuine $\varepsilon\omega$-psh competitor.
Its envelope is nonempty and bounded above by $\Phi_K\le0$, so its
regularization is genuine. Scaling and pullback give the same witnesses
for every approximating and limiting reference.

\section{The blowup, mass and limiting equation}

A genuine $\pi^*\omega$-psh potential descends under a point blowup.
Off the exceptional divisor, add a local downstairs $\omega$-potential.
The resulting psh function is upper bounded near the center because
the upstairs potential is upper bounded near the compact exceptional
fiber. Removable extension gives a canonical downstairs psh potential.
Its pullback agrees with the original off the divisor, hence everywhere
as a canonical representative. Conversely genuine potentials pull back.
This bijection preserves inequalities and envelopes. Scaling therefore
proves the asserted ordinary model identities on $X$.

For completeness, $\Theta$ is big by an explicit K\"ahler current.
Choose coordinates with $q=[1:0:\cdots:0]$ and let
\[
 \psi_q=\log\frac{\sum_{i=1}^n|Z_i|^2}{\sum_{i=0}^n|Z_i|^2}.
\]
The resolved direction map $Q:X\to\CP^{n-1}$ gives, for $0<t<1$,
\[
 \Theta+\ddc(t\pi^*\psi_q)
 =(1-t)\pi^*\omega+tQ^*\omega_{n-1}+t[E_q].
\]
The smooth sum is K\"ahler: $(\pi,Q)$ is the incidence graph embedding
into $\CP^n\times\CP^{n-1}$ and both coefficients are positive.
Thus $X$ is K\"ahler and $\Theta$ is smooth semipositive and big.
The local blowup charts and exceptional normal bundle used here are
described in \cite[VII.12]{Demailly}.

\begin{lemma}\label{lem:mass}
The non-pluripolar reference mass of $\Phi_K$ on $\CP^n$ is
$n!\Vol(K)$; that of $r\pi^*\Phi_K$ is $r^n n!\Vol(K)$.
\end{lemma}
\begin{proof}
Let $M_g$ be the Alexandrov measure of a finite convex function, using
Euclidean volume of subgradient images. With
$z_i=e^{x_i/2+i\vartheta_i}$ and $0\le\vartheta_i<2\pi$, direct
differentiation gives
\[
 \mathrm d^cg=\frac1{2\pi}\sum_i g_i\,\mathrm d\vartheta_i,
 \qquad
 \int_T(\ddc g)^n=n!\int_{\mathbb R^n}\det(D^2g)\,\mathrm dx
\]
for smooth convex $g$. The angular integrals cancel $(2\pi)^n$.
Convolution, local Bedford--Taylor continuity and the Alexandrov
description extend the identity to finite convex $g$
\cite[Sections 2.1--2.2]{BB}.
The numerical coefficient here is computed in our operator and
$x=\log|z|^2$ normalization, rather than imported from another convention.
Local uniform continuity of bounded psh products follows by telescoping
the products against a smooth compactly supported test; each term contains
the potential difference and a mixed product whose local mass is bounded
by repeated cutoff integration by parts \cite[Section 2]{BT}.
The precise locally uniform continuity statement is also
\cite[III.3.6]{Demailly}; the finite convex limits used here are continuous.

For $g=g_K$, all subgradients lie in $K$ by separation and the bound
$g_K\le h_K+\log(n+1)$. Conversely, if $p\in\operatorname{int}K$ and
$p+cB\subset K$, then $g_K(x)-p\cdot x\ge c|x|$.
This function attains a minimum, so $p$ is a subgradient at some point.
Consequently
\[
 \operatorname{int}K\subset\partial g_K(\mathbb R^n)\subset K,
 \qquad M_{g_K}(\mathbb R^n)=\Vol(K).
\]
The boundary of a convex body has zero Euclidean volume. The analytic
complement of $T$ receives no non-pluripolar mass
\cite[Section 1.2]{BEGZ}; thus the computed torus mass is the full mass.
For a blowup, locality gives equality outside the exceptional divisor,
while both non-pluripolar products give that divisor and its image zero
mass. Positive scaling multiplies the top product by $r^n$.
\end{proof}

Apply Lemma~\ref{lem:mass} at each datum, including the limit. The common
simplex gives $m>0$. The smooth closed background is big on compact
K\"ahler $X$, and the reference is an ordinary model. For bounded $F$,
the measure $m e^F\mathrm dV/Z(F)$ has mass exactly $m$ and a bounded
nonnegative density relative to an auxiliary K\"ahler volume.
The precise prescribed-singularity existence theorem
\cite[Theorem A(i)]{DDL} produces $u\sim\phi$ solving that equation.
This proves the limiting solvability; the same argument works for all $j$.
Scalar uniqueness then identifies $c(F)=m/Z(F)$ and $c_j(F)=m_j/Z(F)$.
One may normalize solutions by $\sup_Xu=0$.

The operator conversion in this application is explicit. If the source
uses $D=i\partial\bar\partial$, replace the background by $2\pi\theta$
and the right-hand measure by $(2\pi)^n m e^F\mathrm dV/Z(F)$.
Potentials, psh cones, envelopes and bounded-difference types coincide,
and the matching-mass hypothesis remains exact. Potential uniqueness
is not needed to identify the scalar.

\section{Functional estimates and all support clusters}

The common simplex has an inscribed ball $B(c,s)$ with
$s=\delta/(n+\sqrt n)$. Hausdorff distance gives
$P_j\subset P+h_jB$. Since $sB\subset P-c$, convexity implies
\[
 P+h_jB\subset c+(1+h_j/s)(P-c).
\]
Reverse $P$ and $P_j$ to bound both volume ratios by $(1+h_j/s)^n$.
This proves \eqref{eq:functional-bound}; the simplex volume is
$\delta^n/n!$. The functional difference is independent of the test $G$.

Let $\omega_\ell$ be the uniform-continuity modulus of $\ell$ on $\Sigma$.
Matching maximizers through Hausdorff distance yields
\[
 |g_{P_j}(x)-g_P(x)|\le h_j|x|+\omega_\ell(h_j).
\]
This gives local uniform convergence over $T$ and after pullback.
The common bound $\chi-\delta\log(n+1)\le\Phi_{P_j}\le0$ supplies
an integrable logarithmic lower bound on $X$, including after blowup.
Since $r_j$ is bounded, dominated convergence outside the analytic
coordinate boundary gives the claimed $L^1$ convergence.

For $\nu_F=e^F\mathrm dV/Z(F)$, Jensen's inequality gives
\[
 \log Z(G)-\log Z(F)=\log\int_Xe^{G-F}\mathrm d\nu_F
                \ge\int_X(G-F)\,\mathrm d\nu_F.
\]
Thus $\nu_F$ supports the functional. If another probability $\nu$
supports it, use $G=F+th$ for any continuous real $h$ and both signs
of $t$. Differentiation at zero gives
$\int h\,\mathrm d\nu=\int h\,\mathrm d\nu_F$; hence $\nu=\nu_F$.
If $\norm{F_j-F}_\infty\le\eta$, the density ratio
$\mathrm d\nu_{F_j}/\mathrm d\nu_F$ lies in
$[e^{-2\eta},e^{2\eta}]$, which proves the stated $L^1$ bound.
Theorem~\ref{thm:criterion} follows. Cluster identification has been
proved explicitly, rather than inferred from probability compactness.

\section{Isolated poles of genuinely nonanalytic type}

For $0<a\le1$ let
\[
 P_a=\{p_i\ge0,\ \textstyle\sum_i p_i\le1,
                      \textstyle\sum_i\sqrt{p_i}\ge\sqrt a\}.
\]
The square-root sum is concave, so these are compact convex sets.
Put
\[
 b=\frac{n+1}{2n},\qquad p_{*,i}=\frac bn,
 \qquad\delta=\frac{n-1}{4n}.
\]
Then $b+\delta<1$ and $\sum_i\sqrt{p_{*,i}}=\sqrt{(n+1)/2}>1$.
Adding a point of $\delta\Sigma$ preserves the square-root inequality.
Thus the common simplex is contained in every $P_a$, including $P_1$.

\begin{proposition}\label{prop:nonanalytic}
The ordinary model $\Phi_{P_a}$ has exactly the isolated pole
$p=[1:0:\cdots:0]$ and is not bounded-equivalent to an analytic
singularity. Every generic divisorial Lelong number is zero. Its mass is
\[
 n!\Vol(P_a)=1-\kappa_n a^n,\qquad
 \kappa_n=\frac{2^n n!}{(2n)!}=\frac1{(2n-1)!!}<1.
\]
The same pole assertions hold on $\operatorname{Bl}_q\CP^n$ when $q\ne p$.
\end{proposition}
\begin{proof}
Since all $e_i$ belong to $P_a$,
\[
 \psi_{P_a}\ge
 \log\frac{\max_{1\le i\le n}|Z_i|^2}{\sum_{i=0}^n|Z_i|^2},
\]
which is locally bounded below away from $p$.
Cauchy--Schwarz gives $\sum p_i\ge a/n$ on $P_a$. When all $x_i<0$,
$h_{P_a}(x)\le(a/n)\max_i x_i\to-\infty$ as $z\to p$.
The bounded-equivalence estimate proves the canonical pole assertion.
Away from this single point the potentials are locally bounded, proving
the generic divisorial assertion; pullback at a distinct point adds no pole.

For $\lambda_i>0$ take $|z_i|^2=e^{-R\lambda_i}$.
Writing $y_i=\sqrt{p_i}$ and minimizing by Cauchy--Schwarz gives
\[
 y_i=\frac{\sqrt a}{\lambda_i\sum_k\lambda_k^{-1}},\qquad
 \min_{p\in P_a}\sum_i\lambda_i p_i
       =\frac a{\sum_i\lambda_i^{-1}}=:d_a(\lambda).
\]
The minimizer has $\sum p_i\le a\le1$.
Since $L(-R\lambda)\to0$ and $\Phi_{P_a}-\psi_{P_a}$ is bounded,
$-\Phi_{P_a}/R\to d_a(\lambda)$ along these logarithmic directions.

If this type were $c\log\sum_k|f_k|^2+O(1)$, phase averaging and Taylor
orthogonality would give the positive series
\[
 \sum_{k,\alpha}|c_{k,\alpha}|^2e^{-R\alpha\cdot\lambda}.
\]
Bounded equivalence makes its exponential rate $d_a(\lambda)/c$.
For a convergent positive series the rate is the minimum exponent:
one minimizing term bounds it below, and convergence at a fixed $R_0$
bounds the sum above by a constant times
$e^{-(R-R_0)\min\alpha\cdot\lambda}$.
On a compact neighborhood of a strictly positive $\lambda$, only
finitely many monomials can minimize, since each $\lambda_i$ has a
positive lower bound and one fixed nonzero coefficient bounds the degree.
The resulting rate is locally piecewise linear. But
\[
 d_a(s,1,\ldots,1)=\frac{as}{1+(n-1)s},\qquad
 \frac{\mathrm d^2}{\mathrm ds^2}d_a(s,1,\ldots,1)
   =-\frac{2a(n-1)}{(1+(n-1)s)^3}<0.
\]
This contradicts piecewise linearity on every open interval.

The removed corner $C_a=\{p_i\ge0:\sum_i\sqrt{p_i}<\sqrt a\}$
equals $aC_1$ and lies in $\Sigma$. Substitution $p_i=y_i^2$ gives
\[
 \Vol(C_1)=2^n\int_{y_i\ge0,\,\sum y_i<1}\prod_i y_i\,\mathrm dy
          =\frac{2^n}{(2n)!}.
\]
Iterated monomial integration on a simplex proves the last equality.
Subtracting the corner from $\Sigma$ proves the mass formula.
\end{proof}

For $\theta_r=r\Theta$ and $\phi_{r,a}=r\pi^*\Phi_{P_a}$,
\[
 B_{r,a}(F)=n\log r+\log(1-\kappa_n a^n)-\log Z(F),
 \qquad A_{r,a}(F)=\{\nu_F\}.
\]
Take $a_j=1-1/(j+2)$ and $r_j=1+1/(j+2)$, $j\ge1$.
For $p\in P_{a_j}\setminus P_1$, set $u=(\sum\sqrt{p_i})^2\in[a_j,1)$.
Then $p/u\in P_1$ and $|p/u-p|\le1-u\le1-a_j$.
Consequently $d_H(P_{a_j},P_1)\le1-a_j$ and
\[
 \sup_G|B_{r_j,a_j}(G)-B_{1,1}(G)|
 \le n|\log r_j|+\frac{n\kappa_n}{1-\kappa_n}|1-a_j|\longrightarrow0.
\]
The endpoint $a=1$ is an actual admissible solvable ordinary class of
positive mass $1-\kappa_n$, equal to $2/3$ for $n=2$.
The prototypes have different directional coefficients when $a$ changes;
even the scaled sequence has $r_ja_j=1-(j+2)^{-2}$ strictly varying.
Thus its ordinary types vary, as do its masses.

\section{A smooth degenerate residual with positive mass}

Fix distinct points $p,q\in\CP^n$, take $X=\operatorname{Bl}_q\CP^n$
and $\Theta=\pi^*\omega$. Choose $p=[1:0:\cdots:0]$. For $0<t<1$ set
\[
 \psi_t=t\log\frac{\sum_{i=1}^n|Z_i|^2}{\sum_{i=0}^n|Z_i|^2},
 \qquad K_t=\{p\in\Sigma:\textstyle\sum_i p_i\ge t\}.
\]
The local current potential of $\psi_t$ is
$(1-t)L+t\log\sum e^{x_i}$. Its bounded-distance support function is
$(1-t)h_\Sigma+t\max_i x_i$, the support function of
\[
 (1-t)\Sigma+t\operatorname{conv}(e_1,\ldots,e_n)=K_t.
\]
The log-sum-exp errors are bounded, so
$\psi_t\sim\psi_{K_t}\sim\Phi_{K_t}$.
For $0<t_{\min}\le t\le t_{\max}<1$, choose
$p_{*,i}=t_{\max}/n$ and $\delta=(1-t_{\max})/2$.
This is a common translated simplex in all $K_t$.
For $s\ge t$ in this interval, move a point of $K_t$ toward a fixed
vertex until its coordinate sum reaches $s$. Since $\operatorname{diam}\Sigma=\sqrt2$,
this gives $d_H(K_s,K_t)\le\sqrt2(s-t)/(1-t_{\max})$.
Thus $t_j\to t$ supplies the Hausdorff hypothesis.
Since $n!\Vol(K_t)=1-t^n$, Theorem~\ref{thm:criterion} gives
\[
 \sup_G|B_j(G)-B(G)|
 \le n|\log(r_j/r)|+
       \frac{nt_{\max}^{n-1}}{1-t_{\max}^n}|t_j-t|.
\]
For example, $t_j=1/2+1/[4(j+1)]$ and $r_j=1+1/(j+2)$
have limiting reference mass $1-2^{-n}$ and the explicit support
probabilities $\nu_F$.

Let $Y=\operatorname{Bl}_p\CP^n$, with blowdown $\pi_p$ and resolved
direction map $Q_p:Y\to\CP^{n-1}$. Put
\[
 \beta_t=(1-t)\pi_p^*\omega+tQ_p^*\omega_{n-1}.
\]
This is K\"ahler by the incidence graph embedding.
Resolve the pole on $X$ by $\tau:Z=\operatorname{Bl}_{\pi^{-1}(p)}X\to X$.
The disjoint centers also identify $Z$ with the blowup
$\sigma:Z\to Y$ at the point above $q$. Set $\alpha_t=\sigma^*\beta_t$.
Poincar\'e--Lelong gives
\begin{equation}\label{eq:residual}
 (\pi\circ\tau)^*\omega+\ddc(\psi_t\circ\pi\circ\tau)
                 =\alpha_t+t[E_p].
\end{equation}
There is no $E_q$ coefficient, since the prototype is bounded near $q$.
For the ordinary model, bounded equivalence changes the residual to
$\alpha_t+\ddc w_t$ with a bounded canonical $\alpha_t$-psh $w_t$.
The bounded remainder extends canonically across $E_p$ after adding
local smooth residual potentials. Its value on the divisor is not
inserted into an unbounded product.

\begin{proposition}
For $0<t<1$, the smooth $\alpha_t$ has rank exactly one at every point
of $E_q$, is positive away from $E_q$, and satisfies
$\int_Z\alpha_t^n=1-t^n>0$.
\end{proposition}
\begin{proof}
In a blowup chart the map is
$(s,w_2,\ldots,w_n)\mapsto(s,sw_2,\ldots,sw_n)$.
At $s=0$ its differential has rank one and kills the $n-1$ tangential
directions of the exceptional divisor. Since $\beta_t$ is K\"ahler,
the pullback has exactly the stated rank and is positive elsewhere.

For the mass, write $H=\pi_p^*[\omega]$, $E=[E_p]$.
The current identity gives $[\beta_t]=H-tE$. The restriction of $H$
to $E_p$ is zero, so $\int H^{n-k}E^k=0$ for $1\le k<n$.
The exceptional normal line is $\mathcal O_{\CP^{n-1}}(-1)$
\cite[VII.12.4]{Demailly}; hence $\int E^n=(-1)^{n-1}$.
Expanding gives $\int_Y\beta_t^n=1-t^n$.
The degree-one blowup preserves the top integral, proving the claim.
\end{proof}

Thus there is no positive lower bound for $\alpha_t$ against a K\"ahler
metric, although its top mass is positive. Scaling by $r$ scales the
residual and the exceptional coefficient in \eqref{eq:residual} by $r$.
Solvability is established on $X$ by the model criterion and the precise
existence theorem, without treating the resolved residual as strictly positive.

At $t=1$ the residual is pulled back from $\CP^{n-1}$ and has zero top
integral. Every ordinary-type potential lifts to a bounded residual-psh
remainder. The bounded Bedford--Taylor telescoping identity and Stokes
give the same zero top mass for every such remainder
\cite[Theorem 1.14, with the exceptional set empty]{BEGZ}. Positivity then
forces its top measure to vanish, so no positive-density equation
$c e^F\mathrm dV$ with $c>0$ is solvable in that type. This endpoint is
excluded; the curved-family endpoint $a=1$ has positive mass and is included.

\section{Scope of the conclusion}

The geometric criterion applies to the specified closed projective
backgrounds and fixed blowups. It supplies varying ordinary types with
nondivisorial, genuinely nonanalytic poles and an explicit family of
smooth degenerate positive-mass residuals. It proves the limiting
admissible class, its equation, the functional estimates and every
support-cluster identity. Arbitrary nonclosed Hermitian torsion, moving
modifications and unrestricted singularity degeneration are outside
this theorem. No convergence of solution potentials or weighted rough
adjoint product is asserted. The broader question S08 remains open.

\begin{thebibliography}{9}
\bibitem{BB}
R. J. Berman and B. Berndtsson,
\emph{Real Monge--Amp\`ere equations and K\"ahler--Ricci solitons on toric
log Fano varieties}, arXiv:1207.6128v1 (25 July 2012), Sections 2.1--2.3.
\href{https://arxiv.org/abs/1207.6128v1}{Official primary version}.
\bibitem{DDL}
T. Darvas, E. Di Nezza and C. H. Lu,
\emph{Log-concavity of volume and complex Monge--Amp\`ere equations with
prescribed singularity}, arXiv:1807.00276v3 (12 March 2020).
\href{https://arxiv.org/abs/1807.00276v3}{Official primary version}.
\bibitem{Demailly}
J.-P. Demailly, \emph{Complex Analytic and Differential Geometry},
Chapters I, III and VII.
\href{https://www-fourier.univ-grenoble-alpes.fr/~demailly/manuscripts/agbook.pdf}
{Author's manuscript}.
\bibitem{BT}
E. Bedford and B. A. Taylor,
\emph{A new capacity for plurisubharmonic functions},
Acta Math. \textbf{149} (1982), 1--40.
\bibitem{BEGZ}
S. Boucksom, P. Eyssidieux, V. Guedj and A. Zeriahi,
\emph{Monge--Amp\`ere equations in big cohomology classes},
Acta Math. \textbf{205} (2010), 199--262.
\end{thebibliography}
\end{document}
